\hypertarget{cont5:paper-ix-continuous-p-adic-parameter-charts}{%
\subsection{Paper IX: continuous p-adic parameter
charts}\label{cont5:paper-ix-continuous-p-adic-parameter-charts}}

\textbf{Komorebi P. --- 8 September 2026.}

\textbf{Fresh verification: PASS --- 40 checks, 244 explicit declaration
audits.}

\hypertarget{cont5:main-statement}{%
\subsubsection{Main statement}\label{cont5:main-statement}}

Let \(p\) be prime, \(L\ge1\), \(e,q\ge0\), and \[
F\in\mathbb Z[X_1,\ldots,X_{e+q}]^e,\qquad
x\in\mathbb Z^{e+q},\qquad F(x)\equiv0\pmod{p^L}.
\] Assume the actual reduced Jacobian \[
A_x=DF(x)\bmod p:\mathbb F_p^{e+q}\longrightarrow\mathbb F_p^e
\] is surjective. All vector congruences are coordinatewise.

There is a coordinate permutation \(\sigma\), with
\((\sigma z)_i=z_{\sigma(i)}\). Write \(\sigma z=(y,t)\), using \(e\)
solved coordinates and \(q\) free coordinates. Define \[
(y_0,t_0)=\sigma x,\qquad
B=x+p^L\mathbb Z_p^{e+q},\qquad
T=t_0+p^L\mathbb Z_p^q,\qquad
\mathcal R=\{z\in B:F(z)=0\}.
\] There is a homeomorphism \[
h:T\xrightarrow{\ \sim\ }\mathcal R,\qquad
h^{-1}(z)=\operatorname{free}(\sigma z).
\] Equivalently, a continuous function
\(\phi:T\to y_0+p^L\mathbb Z_p^e\) gives the entire root locus through
\(\sigma h(t)=(\phi(t),t)\). Every parameter has exactly one
corresponding root in \(B\). For every \(N\ge0\), \[
\boxed{h(t)\equiv h(u)\pmod{p^N}\iff t\equiv u\pmod{p^N}.}
\]

The coordinate choice is a conclusion. A suitable minor and a
preexisting chart are not hypotheses. The theorem covers the original
root locus after undoing the coordinate reordering. Since \(x\) is an
approximate root, it does not assert \(\phi(t_0)=y_0\) unless \(F(x)=0\)
exactly. The cases \(p=2\), \(L=1\), \(e=0\), and \(q=0\) are included.

\hypertarget{cont5:supporting-theorems}{%
\subsubsection{Supporting theorems}\label{cont5:supporting-theorems}}

For a rectangular integer polynomial system with a surjective reduced
Jacobian at \(x\), and \(L\ge1\), \[
F(x+p^L\mathbb Z_p^d)=F(x)+p^L\mathbb Z_p^e.
\] This covers arbitrary p-adic output targets. No lower-root hypothesis
is needed.

When \(d=e\), the actual restricted polynomial map is a homeomorphism.
Moreover, for arbitrary \(a,b\in x+p\mathbb Z_p^d\) and every \(N\ge0\),
\[
F(a)\equiv F(b)\pmod{p^N}\iff a\equiv b\pmod{p^N}.
\] Its inverse therefore preserves every congruence depth.

\hypertarget{cont5:proof-structure}{%
\subsubsection{Proof structure}\label{cont5:proof-structure}}

\begin{enumerate}
\def\labelenumi{\arabic{enumi}.}
\tightlist
\item
  The actual integer Taylor formula gives one additional input digit
  from one additional output digit when the reduced square Jacobian is
  surjective. Its finite domain and codomain have the same size, so it
  is injective. Induction gives every integer precision.
\item
  Canonical integer truncations of arbitrary p-adic vectors preserve
  evaluation modulo each power of \(p\). Applying the integer result to
  truncations proves the p-adic congruence equivalence.
\item
  Subtracting an integer target from the constant terms leaves the
  actual Jacobian unchanged. The previous root-lifting theorem supplies
  a preimage. The image of the compact input ball is closed and contains
  the convergent truncations of every admissible p-adic target, hence
  contains the target. No compatibility of arbitrarily chosen preimages
  is assumed.
\item
  In the square case, the resulting continuous bijection from a compact
  ball to a Hausdorff ball is a homeomorphism. Its forward map is actual
  polynomial evaluation, and its inverse solves the output equation.
\item
  For a specified surjective solved-variable block, apply this inverse
  to the actual augmented polynomial map \[
  \Psi(y,t)=(F(y,t),t).
  \] Its derivative has an identity free-coordinate block. The inverse
  at \((0,t)\) gives the root. The lower-root congruence places this
  target in the output ball. Injectivity proves uniqueness and full
  graph coverage. Projection identifies the inverse chart and reflects
  congruences.
\item
  The columns of a surjective finite-dimensional linear map contain a
  basis. Select those columns and extend their injection to a coordinate
  permutation. Polynomial renaming by \(\sigma^{-1}\) transports actual
  evaluations and partial derivatives. The fixed-coordinate chart then
  gives the main theorem on the entire original root locus.
\end{enumerate}

\hypertarget{cont5:scope}{%
\subsubsection{Scope}\label{cont5:scope}}

Coefficients are integers; roots and parameters are arbitrary p-adic
points. This supplies continuous polynomial parameter charts.
Analyticity, differentiability, arbitrary p-adic coefficient interfaces,
weighted lifting for nonunit Jacobians, and first-order root-motion
estimates are separate.

Full Jacobian surjectivity need not allow an arbitrary prescribed
parameter choice. At the dyadic origin, \(F(X,T)=X^2-T\) has a
surjective full Jacobian but zero \(X\)-derivative. Prescribing \(T=2\)
gives no root even modulo \(4\). Suitable coordinates instead solve for
\(T\).

\hypertarget{cont5:verification-and-provenance}{%
\subsubsection{Verification and
provenance}\label{cont5:verification-and-provenance}}

The eight new mathematical modules contain 75 declarations: 60 theorems
and 15 definitions. The active polynomial dependency set contains 25
mathematical modules. Counts including inherited support and theorem
contracts are recorded separately in the final report.

The exact universal statements are Lean theorems, not conclusions from
sampled numerical tests. The fresh gate rebuilds project sources with
warnings treated as errors, checks expanded theorem contracts, audits
explicit declarations, and checks rejected invalid extensions. Only
propext, Classical.choice, and Quot.sound are permitted as axiom
dependencies.

Pinned runtime: Lean 4.34.0-rc2, Mathlib revision
\mcode{aeb72f1e9981b6d35b426b16dca7759be3e7ba06}.

See \nolinkurl{inventory.json}, \nolinkurl{verification/polynomial_charts_verification.json},
and \nolinkurl{INDEPENDENT_PROOF_AUDIT.md} for actual evidence.
RECOVERY\_RECORD.md explains workspace maintenance, the intact previous
release, and reconstruction of the unpublished chart sources. Lost logs
are not represented as recovered.
